\documentclass[11pt,a4paper]{article}
\usepackage{amsmath,amssymb,amsthm}
\usepackage{bm}
\usepackage{geometry}
\geometry{margin=2.5cm}
\usepackage{booktabs}
\usepackage[pdfusetitle,hidelinks]{hyperref}

\newcommand{\norm}[1]{\left\lVert #1 \right\rVert}
\newcommand{\Fro}{\mathrm{F}}
\newcommand{\diag}{\operatorname{diag}}
\newcommand{\Reop}{\operatorname{Re}}
\newcommand{\Sigmat}{\bm{\Sigma}}

\title{Fair Reconstruction Error for SPDMD\\with and without Mean Subtraction}
\author{Internal note — Project 1}
\date{}
\begin{document}
\maketitle

%======================================================================
\section{SPDMD Reconstruction in Reduced Space}
%======================================================================

The snapshot matrix after preprocessing (optional mean subtraction and/or linear detrending) is
\begin{equation}
\bm{D} = \begin{bmatrix} \bm{x}_1 & \bm{x}_2 & \cdots & \bm{x}_{m-1} \end{bmatrix} \in \mathbb{R}^{n \times (m-1)},
\end{equation}
where $n$ is the spatial dimension and $m$ the total number of snapshots.

The economy SVD of $\bm{D}$ (truncated to rank $r$) gives
\begin{equation}\label{eq:svd}
\bm{D} \approx \bm{U}_r\,\Sigmat_r\,\bm{V}_r^\top,
\end{equation}
where $\bm{U}_r \in \mathbb{R}^{n\times r}$, $\Sigmat_r = \diag(\sigma_1,\ldots,\sigma_r) \in \mathbb{R}^{r\times r}$, $\bm{V}_r \in \mathbb{R}^{(m-1)\times r}$.
In the code, \texttt{SVD\_S} stores $\Sigmat_r$ and \texttt{SVD\_VT} stores $\bm{V}_r^\top$.

\subsection{Reduced data matrix \texorpdfstring{$\bm{G}$}{G}}

The data projected into the $r$-dimensional reduced space is
\begin{equation}\label{eq:G}
\bm{G} = \Sigmat_r\,\bm{V}_r^\top \;\in\; \mathbb{R}^{r \times (m-1)}.
\end{equation}
This is computed in \texttt{dmd.f95} as
\begin{verbatim}
G = matmul(transpose(SVD_S), transpose(SVD_VT))
\end{verbatim}
Note that $\bm{G} = \bm{U}_r^\top \bm{D}$, i.e.\ the projection of $\bm{D}$ onto its left singular vectors.

\subsection{DMD eigenvalues and the Vandermonde matrix}

After computing the reduced dynamics matrix $\tilde{\bm{A}} = \bm{U}_r^\top \bm{D}''\bm{V}_r\Sigmat_r^{-1}$ and solving its eigenvalue problem $\tilde{\bm{A}}\bm{W} = \bm{W}\bm{\Lambda}$, the DMD eigenvalues $\mu_1,\ldots,\mu_r$ define the Vandermonde matrix
\begin{equation}\label{eq:Vand}
\bm{C} \;\in\; \mathbb{C}^{r \times (m-1)}, \qquad C_{j,k} = \mu_j^{\,k-1}, \quad j=1,\ldots,r,\;\; k=1,\ldots,m-1.
\end{equation}

\subsection{SPDMD sparse amplitudes}

The sparsity-promoting optimisation (Jovanovi\'c et al.\ 2014) solves, for each penalty $\gamma$,
\begin{equation}
\min_{\bm{b}} \;\underbrace{\norm{\bm{G} - \bm{W}\,\diag(\bm{b})\,\bm{C}}_\Fro^2}_{J_1} + \gamma\,\norm{\bm{b}}_1,
\end{equation}
where $\bm{W} = \bm{V}_r$ (the DMD eigenvectors in the reduced basis, stored as \texttt{Vr}) and $\bm{b}\in\mathbb{C}^r$ are the mode amplitudes.  The matrices
\begin{align}
\bm{P} &= \left(\bm{W}^H \bm{W}\right) \circ \overline{\left(\bm{C}\,\bm{C}^H\right)} \;\in\;\mathbb{C}^{r\times r}, \label{eq:P}\\
\bm{q} &= \diag\!\left(\bm{W}^H\,\bm{G}\,\bm{C}^H\right) \;\in\;\mathbb{C}^{r}, \label{eq:q}
\end{align}
(where $\circ$ denotes Hadamard product) reduce the least-squares term to
\begin{equation}
J_1 = \bm{b}^H \bm{P}\,\bm{b} - 2\,\Reop\!\left(\bm{q}^H \bm{b}\right) + \norm{\bm{G}}_\Fro^2.
\end{equation}
The ADMM solver (\texttt{amp.f95}, module \texttt{Sparity.f95}) produces sparse amplitudes $\bm{b}_\gamma$ for each $\gamma$ value.

%======================================================================
\section{Reconstruction Error}
%======================================================================

Given sparse amplitudes $\bm{b}_\gamma$ with \texttt{nzval} nonzero entries, the reduced-space reconstruction is
\begin{equation}\label{eq:Xdmd}
\hat{\bm{G}}_\gamma = \Reop\!\left(\bm{W}\,\diag(\bm{b}_\gamma)\,\bm{C}\right) \;\in\; \mathbb{R}^{r \times (m-1)}.
\end{equation}
This is computed in \texttt{reconstruction\_opt\_error.f95} as:
\begin{verbatim}
B(j,j) = B_p(j, r4)       ! diagonal matrix of sparse amplitudes
xdmd = real(matmul(matmul(Vr, B), C))
\end{verbatim}

The reconstruction error is
\begin{equation}\label{eq:eps}
\boxed{\epsilon(\gamma) = \frac{\norm{\bm{G} - \hat{\bm{G}}_\gamma}_\Fro^2}{\norm{\Sigmat_r}_\Fro^2} \times 100\%.}
\end{equation}

The denominator $\norm{\Sigmat_r}_\Fro^2 = \sum_{i=1}^{r}\sigma_i^2$ equals the total energy captured by the rank-$r$ truncation.  In the code:
\begin{verbatim}
F_norm_orgi = sum(SVD_S**2)       ! = ||Sigma_r||_F^2
F_norm_XDMD = sum((G_matrix - xdmd)**2)  ! = ||G - G_hat||_F^2
error = (F_norm_XDMD / F_norm_orgi) * 100
\end{verbatim}


%======================================================================
\section{The Denominator Difference: \texttt{mean0} vs.\ \texttt{mean1}}
\label{sec:denom}
%======================================================================

The crucial observation is that \textbf{the SVD is performed on different data} depending on whether the temporal mean is subtracted:

\begin{itemize}
\item \textbf{\texttt{mean0}} (no mean subtraction): $\bm{D}^{(0)} = \bm{X}$ (raw snapshots $\to$ SVD gives $\sigma_i^{(0)}$).
\item \textbf{\texttt{mean1}} (mean subtracted): $\bm{D}^{(1)} = \bm{X} - \bar{\bm{x}}\,\bm{1}^\top$ (fluctuations $\to$ SVD gives $\sigma_i^{(1)}$).
\end{itemize}

Since $\bm{D}^{(0)}$ contains the mean field, its singular values are larger:
\begin{equation}
\norm{\Sigmat_r^{(0)}}_\Fro^2 \;>\; \norm{\Sigmat_r^{(1)}}_\Fro^2.
\end{equation}

The respective reconstruction errors are
\begin{align}
\epsilon^{(0)} &= \frac{\norm{\bm{G}^{(0)} - \hat{\bm{G}}^{(0)}}_\Fro^2}{\norm{\Sigmat_r^{(0)}}_\Fro^2} \times 100\%,\label{eq:eps0}\\[6pt]
\epsilon^{(1)} &= \frac{\norm{\bm{G}^{(1)} - \hat{\bm{G}}^{(1)}}_\Fro^2}{\norm{\Sigmat_r^{(1)}}_\Fro^2} \times 100\%.\label{eq:eps1}
\end{align}
These use \textbf{different denominators}, so comparing $\epsilon^{(0)}$ to $\epsilon^{(1)}$ directly is not meaningful.

%======================================================================
\section{Fair Error Metric}
\label{sec:fair}
%======================================================================

To compare reconstruction quality on the same footing, we normalise $\epsilon^{(1)}$ so that both errors use the \emph{same denominator} — the total energy including the mean:

\begin{equation}\label{eq:alpha}
\alpha \;=\; \frac{\norm{\Sigmat_r^{(1)}}_\Fro^2}{\norm{\Sigmat_r^{(0)}}_\Fro^2}
\;=\; \frac{\sum_{i=1}^{r} \left(\sigma_i^{(1)}\right)^2}{\sum_{i=1}^{r} \left(\sigma_i^{(0)}\right)^2}.
\end{equation}

This is the \textbf{fluctuation-to-total energy ratio}.  Then the fair reconstruction error for \texttt{mean1} is
\begin{equation}\label{eq:fair}
\boxed{\epsilon_{\mathrm{fair}}^{(1)} \;=\; \alpha \;\cdot\; \epsilon^{(1)} \;=\; \frac{\norm{\bm{G}^{(1)} - \hat{\bm{G}}^{(1)}}_\Fro^2}{\norm{\Sigmat_r^{(0)}}_\Fro^2} \times 100\%.}
\end{equation}

For \texttt{mean0}, the fair error equals the original error:
\begin{equation}
\epsilon_{\mathrm{fair}}^{(0)} \;=\; \epsilon^{(0)}.
\end{equation}

\subsection{Derivation}

The scaling is simply a change of denominator:
\begin{align}
\epsilon_{\mathrm{fair}}^{(1)} 
&= \alpha \cdot \epsilon^{(1)} \nonumber\\
&= \frac{\norm{\Sigmat_r^{(1)}}_\Fro^2}{\norm{\Sigmat_r^{(0)}}_\Fro^2} \;\cdot\;
   \frac{\norm{\bm{G}^{(1)} - \hat{\bm{G}}^{(1)}}_\Fro^2}{\norm{\Sigmat_r^{(1)}}_\Fro^2} \times 100\% \nonumber\\
&= \frac{\norm{\bm{G}^{(1)} - \hat{\bm{G}}^{(1)}}_\Fro^2}{\norm{\Sigmat_r^{(0)}}_\Fro^2} \times 100\%.
\end{align}
The $\norm{\Sigmat_r^{(1)}}_\Fro^2$ cancels, leaving the reconstruction residual of the fluctuation field normalised by the total energy.  Both $\epsilon_{\mathrm{fair}}^{(0)}$ and $\epsilon_{\mathrm{fair}}^{(1)}$ are now fractions of the same quantity.

\subsection{Exactness of \texorpdfstring{$\alpha$}{alpha}: no truncation error}

A potential concern is whether $\alpha$ is affected by SVD truncation.  In general, if the SVD were truncated to some rank $r < m-1$, the sums $\sum_{i=1}^{r}\sigma_i^2$ would discard the tail of the spectrum and $\alpha$ would only be an approximation.

In our implementation, however, the SVD rank equals $m-1$ in all cases (verified across all 39 HDF5 result files).  Since $\bm{D}\in\mathbb{R}^{n\times(m-1)}$ with $n \gg m-1$, the matrix has at most $m-1$ nonzero singular values, and the economy SVD with $r = m-1$ captures \textbf{all of them}.  Therefore
\begin{equation}
\norm{\Sigmat_r}_\Fro^2 = \sum_{i=1}^{m-1}\sigma_i^2 = \norm{\bm{D}}_\Fro^2,
\end{equation}
and the energy ratio is \emph{exact}:
\begin{equation}
\alpha = \frac{\norm{\bm{D}^{(1)}}_\Fro^2}{\norm{\bm{D}^{(0)}}_\Fro^2} = \frac{\text{total fluctuation energy}}{\text{total energy (incl.\ mean)}}.
\end{equation}
No singular values are discarded, so $\alpha$ involves no truncation approximation whatsoever.

Note that the sparsity parameter $\gamma$ only affects the \emph{numerator} of the error, i.e.\ how many DMD modes are used to approximate $\bm{G}$.  The \emph{denominator} $\norm{\Sigmat_r}_\Fro^2$ always uses all $m-1$ singular values regardless of how many modes are retained. Alpha simply re-scales this denominator to be the same for both \texttt{mean0} and \texttt{mean1}.

\subsection{Physical interpretation of \texorpdfstring{$\alpha$}{alpha}}

\begin{itemize}
\item $\alpha \approx 1$: The mean carries little energy; the data is fluctuation-dominated.  Fair and raw errors are similar.  \emph{Example:}\ Eulerian gas fraction, where $\alpha \approx 0.82\text{--}0.92$.
\item $\alpha \ll 1$: The mean dominates the total energy.  The raw \texttt{mean1} error $\epsilon^{(1)}$ is inflated (normalised by small fluctuation energy), making \texttt{mean0} \emph{appear} much better.  \emph{Example:}\ Moving-Eulerian gas fraction, where $\alpha \approx 0.03\text{--}0.31$; Lagrangian gas fraction $\alpha \approx 0.03\text{--}0.08$.
\end{itemize}

%======================================================================
\section{Why the Na\"ive Comparison is Misleading}
%======================================================================

Consider Lagrangian gas fraction, Case C$_1$ (Re20\,Bo20):
\begin{center}
\begin{tabular}{@{}lccc@{}}
\toprule
Variant & $\epsilon$ (raw) & $\alpha$ & $\epsilon_{\mathrm{fair}}$ \\
\midrule
\texttt{mean0} & $0.04\%$ & $0.026$ & $0.04\%$ \\
\texttt{mean1} & $4.33\%$ & $0.026$ & $0.11\%$ \\
\bottomrule
\end{tabular}
\end{center}

Na\"ively, \texttt{mean0} appears $\sim\!100\times$ better ($0.04\%$ vs.\ $4.33\%$).  After fair correction, \texttt{mean0} is only $\sim\!3\times$ better ($0.04\%$ vs.\ $0.11\%$) — and this residual difference is on fluctuation energy that is only $2.6\%$ of the total, i.e.\ both variants reconstruct $>99.9\%$ of the total field energy.

The na\"ive comparison is misleading because the \texttt{mean0} denominator $\norm{\Sigmat_r^{(0)}}_\Fro^2$ includes the mean energy (which is trivially reproduced), inflating the denominator and shrinking the percentage.

%======================================================================
\section{Summary}
%======================================================================

\begin{enumerate}
\item The standard SPDMD error metric $\epsilon = \norm{\bm{G}-\hat{\bm{G}}}_\Fro^2 / \norm{\Sigmat_r}_\Fro^2 \times 100\%$ uses $\Sigmat_r$ from whichever data was decomposed.
\item When comparing \texttt{mean0} and \texttt{mean1}, the denominators differ: \texttt{mean0} includes mean energy, \texttt{mean1} does not.
\item The fair error $\epsilon_{\mathrm{fair}}^{(1)} = \alpha\,\epsilon^{(1)}$ normalises both by the same total energy.
\item After fair correction, \texttt{mean0} and \texttt{mean1} produce \textbf{comparable} reconstruction quality for non-Eulerian frames.  For Eulerian frames, \texttt{mean1} is genuinely much better.
\item Mean subtraction is numerically \textbf{essential} for SPDMD in non-Eulerian frames (prevents Vandermonde overflow and ADMM failure) and costs nothing in reconstruction accuracy.
\end{enumerate}


%======================================================================
%======================================================================
\newpage
\section{Manuscript--Code Consistency Check}
\label{sec:consistency}
%======================================================================
%======================================================================

This section verifies that the theoretical formulation in the manuscript (Section~2, ``Data-driven modal decomposition methods'') is consistent with the Fortran implementation, and provides a complete variable mapping.

\subsection{Variable mapping: manuscript \texorpdfstring{$\longleftrightarrow$}{<->} code}

\begin{center}
\renewcommand{\arraystretch}{1.25}
\begin{tabular}{@{}lll p{6.2cm}@{}}
\toprule
\textbf{Manuscript} & \textbf{Code variable} & \textbf{File} & \textbf{Description} \\
\midrule
$\mathcal{D}'$ & \texttt{u\_stacked\_full(:,1:m-1)} & \texttt{dmd.f95} & First $m{-}1$ snapshots after preprocessing (mean/trend subtracted) \\
$\mathcal{D}''$ & Built from snapshots $2..m$ & \texttt{dmd.f95} & Shifted snapshot matrix \\
$U_r$ & \texttt{SVD\_U} $(n \times r)$ & \texttt{dmd.f95} & Left singular vectors \\
$\Sigma_r$ & \texttt{SVD\_S} $(r \times r)$ & \texttt{dmd.f95} & Diagonal matrix of singular values \\
$W_r$ (manuscript) & \texttt{SVD\_VT} $(m{-}1 \times r)$ & \texttt{dmd.f95} & Right singular vectors $V_r$ \newline \textbf{Note:} code stores $V_r$, not $V_r^T$ despite the name \\
$\tilde{A}=U_r^T\mathcal{D}''W_r\Sigma_r^{-1}$ & \texttt{A} $\to$ \texttt{A\_comp} $(r\times r)$ & \texttt{dmd.f95} & Reduced dynamics matrix \\
$\gamma_j$ (manuscript Eq.\,8) & \texttt{Mu2}, \texttt{Mu} & \texttt{dmd.f95} & DMD eigenvalues $\mu_j$ \\
$\tilde{v}_j$ & \texttt{Vr} $(r \times r)$ from \texttt{zgeev} & \texttt{dmd.f95} & Eigenvectors of $\tilde{A}$ \\
$\varphi_j = U_r \tilde{v}_j$ & \texttt{Phi} $(n_c \times r)$ & \texttt{recon\_opt\_memory.f95} & Full-space DMD modes (only in full reconstruction) \\
$V_\gamma$ (Vandermonde) & \texttt{Vand}, \texttt{C} $(r \times m{-}1)$ & \texttt{dmd.f95}, \texttt{recon\_opt\_error.f95} & $C_{j,k} = \mu_j^{k-1}$ \\
$B_b = \diag(\bm{b})$ & \texttt{B} $= \diag($\texttt{B\_p(:,}\,$\gamma$\texttt{\_idx))}) & \texttt{recon\_opt\_error.f95} & Diagonal sparse amplitudes \\
--- & \texttt{G} $= \Sigma_r V_r^T$ $(r \times m{-}1)$ & \texttt{dmd.f95} & Reduced data matrix (not in manuscript) \\
\bottomrule
\end{tabular}
\end{center}

\subsection{Notation issue: eigenvalue symbol clash}
\label{sec:notation}

The manuscript uses $\gamma_j$ for the eigenvalues of $\tilde{A}$ (Eq.\,8).  However, in the SPDMD literature (Jovanovi\'c et al.\ 2014) and in the code, $\gamma$ is the \textbf{sparsity penalty parameter}, while the eigenvalues are denoted $\mu_j$.  This creates a notational conflict: the SPDMD optimisation is
\[
\min_{\bm{b}} \norm{\bm{G} - \bm{W}\diag(\bm{b})\bm{C}}_\Fro^2 + \gamma \norm{\bm{b}}_1,
\]
where $\gamma$ is the penalty weight and $\mu_j$ are the eigenvalues that build the Vandermonde matrix $\bm{C}$.  \textbf{Recommendation:} Use $\mu_j$ for DMD eigenvalues throughout the manuscript to avoid ambiguity with the SPDMD penalty $\gamma$.

\subsection{Derivation: from manuscript (full-space) to code (reduced-space)}
\label{sec:derivation}

The manuscript states the SPDMD optimisation problem as (paraphrased):
\begin{equation}\label{eq:full_space}
\min_{\bm{b}} \;\norm{\mathcal{D}' - \bm{\varphi}\, B_b\, V_\gamma}_\Fro^2.
\end{equation}
This is in \textbf{full space} ($n_s \times (m{-}1)$).  The code instead works in the \textbf{reduced $r$-dimensional space}.  We now show these are mathematically equivalent.

\paragraph{Step 1: Substitute $\bm{\varphi} = U_r \bm{W}$.}

From the manuscript (Eq.\,9), the DMD spatial modes are $\varphi_j = U_r \tilde{v}_j$, so $\bm{\varphi} = U_r \bm{W}$, where $\bm{W} \in \mathbb{C}^{r \times r}$ holds the eigenvectors of $\tilde{A}$ as columns (code variable \texttt{Vr}).  Thus
\begin{equation}
\mathcal{D}' - \bm{\varphi}\, B_b\, V_\gamma = \mathcal{D}' - U_r\, \bm{W}\, B_b\, \bm{C}.
\end{equation}

\paragraph{Step 2: Express $\mathcal{D}'$ in the SVD basis.}

From the truncated SVD (Eq.\,7): $\mathcal{D}' = U_r \Sigmat_r V_r^T + \bm{E}_{\mathrm{trunc}}$, where $\bm{E}_{\mathrm{trunc}}$ is the truncation residual.  Since $r = m{-}1$ in our case, $\bm{E}_{\mathrm{trunc}} = \bm{0}$ and we have exactly:
\begin{equation}
\mathcal{D}' = U_r \underbrace{\Sigmat_r V_r^T}_{\displaystyle = \bm{G}}.
\end{equation}

\paragraph{Step 3: Use orthonormality of $U_r$.}

Since $U_r^T U_r = I_r$, the Frobenius norm is preserved under left multiplication by $U_r$:
\begin{align}
\norm{\mathcal{D}' - U_r \bm{W} B_b \bm{C}}_\Fro^2
&= \norm{U_r \bm{G} - U_r \bm{W} B_b \bm{C}}_\Fro^2 \nonumber\\
&= \norm{U_r \left(\bm{G} - \bm{W} B_b \bm{C}\right)}_\Fro^2 \nonumber\\
&= \norm{\bm{G} - \bm{W}\, B_b\, \bm{C}}_\Fro^2. \label{eq:reduced_equiv}
\end{align}
The last step holds because for any matrix $\bm{X}$ and orthonormal $U_r$:
\[
\norm{U_r \bm{X}}_\Fro^2 = \mathrm{tr}\!\left(\bm{X}^T U_r^T U_r \bm{X}\right) = \mathrm{tr}\!\left(\bm{X}^T \bm{X}\right) = \norm{\bm{X}}_\Fro^2.
\]

\paragraph{Conclusion.}  The full-space problem~\eqref{eq:full_space} is \textbf{exactly equivalent} to the reduced-space problem solved by the code:
\begin{equation}\label{eq:reduced_space}
\boxed{\min_{\bm{b}} \;\norm{\bm{G} - \bm{W}\,\diag(\bm{b})\,\bm{C}}_\Fro^2
\;=\; \min_{\bm{b}} \;\norm{\mathcal{D}' - \bm{\varphi}\,\diag(\bm{b})\,V_\gamma}_\Fro^2.}
\end{equation}
This equivalence holds exactly because $r = m{-}1$ (no SVD truncation).  Even if $r < m{-}1$, the optimisation would still be equivalent (the truncation error $\norm{\bm{E}_{\mathrm{trunc}}}_\Fro^2$ is constant w.r.t.\ $\bm{b}$ and drops out of the $\arg\min$), though the error \emph{values} would differ by that constant.


\subsection{Reconstruction error: code formula}
\label{sec:error_code}

The reconstruction in reduced space (code \texttt{reconstruction\_opt\_error.f95}):
\begin{align}
\hat{\bm{G}} &= \Reop\!\left(\bm{W}\,\diag(\bm{b}_\gamma)\,\bm{C}\right), \\[4pt]
\epsilon &= \frac{\norm{\bm{G} - \hat{\bm{G}}}_\Fro^2}{\norm{\Sigmat_r}_\Fro^2} \times 100\%.
\end{align}

Using the equivalence~\eqref{eq:reduced_equiv}, this is identical to the full-space relative error:
\begin{equation}
\epsilon = \frac{\norm{\mathcal{D}' - \hat{\mathcal{D}}'}_\Fro^2}{\norm{\mathcal{D}'}_\Fro^2} \times 100\%,
\end{equation}
where $\hat{\mathcal{D}}' = U_r \hat{\bm{G}} = \bm{\varphi}\,\diag(\bm{b}_\gamma)\,V_\gamma$ and the denominator uses
\[
\norm{\mathcal{D}'}_\Fro^2 = \norm{U_r \Sigmat_r V_r^T}_\Fro^2 = \norm{\Sigmat_r}_\Fro^2 = \sum_{i=1}^{r} \sigma_i^2.
\]

\medskip\noindent
\textbf{Code implementation} (lines 288--298 of \texttt{reconstruction\_opt\_error.f95}):
\begin{verbatim}
xdmd = real(matmul(matmul(Vr, B), C))       ! G_hat = Re(W * diag(b) * C)

summ = 0.0_dp
do i = 1, r1
    do j = 1, m-1
        summ = summ + (G_matrix(i,j) - xdmd(i,j))**2
    end do
end do
F_norm_XDMD = summ                           ! ||G - G_hat||_F^2

error = (F_norm_XDMD / F_norm_orgi) * 100    ! F_norm_orgi = sum(SVD_S**2)
\end{verbatim}

\subsection{The SPDMD matrices \texorpdfstring{$\bm{P}$}{P} and \texorpdfstring{$\bm{q}$}{q}: expanding the quadratic}

The manuscript states that ``using SVD properties, the definition of matrix $\varphi$, and some algebra, the result is then in a quadratic form that is solved using ADMM.''  Here we show this algebra explicitly, connecting to Jovanovi\'c et al.\ (2014).

Expanding the reduced-space objective:
\begin{align}
J(\bm{b}) &= \norm{\bm{G} - \bm{W}\diag(\bm{b})\bm{C}}_\Fro^2 \nonumber\\
&= \mathrm{tr}\!\left[(\bm{G} - \bm{W}\diag(\bm{b})\bm{C})^H (\bm{G} - \bm{W}\diag(\bm{b})\bm{C})\right] \nonumber\\
&= \norm{\bm{G}}_\Fro^2 - 2\,\Reop(\bm{q}^H \bm{b}) + \bm{b}^H \bm{P}\, \bm{b}, \label{eq:quadratic}
\end{align}
where
\begin{align}
P_{ij} &= (\bm{W}^H \bm{W})_{ij}\;\overline{(\bm{C}\bm{C}^H)_{ij}}, \label{eq:P_def}\\
q_i &= \left[\diag\!\left(\overline{\bm{C}\,\bm{G}^T\,\bm{W}}\right)\right]_i. \label{eq:q_def}
\end{align}

These are the Hadamard (element-wise) products from Jovanovi\'c et al.\ (2014, Eq.\,11--12).
In \texttt{dmd.f95} (lines $\sim$1033--1044):
\begin{verbatim}
P = matmul(conjg(transpose(Vr)), Vr) &          ! W^H * W
    * conjg(matmul(Vand, conjg(transpose(Vand)))) ! .* conj(C * C^H)

G = matmul(transpose(SVD_S), transpose(SVD_VT))  ! Sigma_r * V_r^T

q_diag = conjg(matmul(matmul(Vand, transpose(G)), Vr))  ! conj(C * G^T * W)
do i = 1, rank
    q(i) = q_diag(i,i)                           ! take diagonal
end do
\end{verbatim}

\noindent This matches Eqs.~\eqref{eq:P_def}--\eqref{eq:q_def} exactly.  The ADMM solver in \texttt{Sparity.f95} then minimises $J(\bm{b}) + \gamma\norm{\bm{b}}_1$ over $\bm{b}$.


\subsection{Verdict: manuscript vs.\ code}
\label{sec:verdict}

\begin{center}
\renewcommand{\arraystretch}{1.2}
\begin{tabular}{@{}p{4cm}cc p{5.5cm}@{}}
\toprule
\textbf{Item} & \textbf{Match?} & & \textbf{Comment} \\
\midrule
SVD of $\mathcal{D}'$ & \checkmark & & Manuscript Eq.\,7 $=$ code \texttt{dgesvd}. \\
$\tilde{A} = U_r^T \mathcal{D}'' W_r \Sigma_r^{-1}$ & \checkmark & & Manuscript Eq.\,8 $=$ code \texttt{A = temp\_two * Inv\_Sigma}. SVD variable naming mismatch: manuscript calls right singular vectors $W_r$, code names them \texttt{SVD\_VT} (actually stores $V_r$).  \\
Eigenvalue problem & \checkmark & & Manuscript Eq.\,8 $=$ code \texttt{zgeev(A\_comp,$\ldots$)}. \\
Spatial modes $\varphi_j = U_r \tilde{v}_j$ & \checkmark & & Manuscript Eq.\,9 $=$ code \texttt{Phi}. \\
SPDMD optimisation & \checkmark & & Manuscript (full-space) $\equiv$ code (reduced-space) via $U_r$ orthonormality. \\
Reconstruction & \checkmark & & $\hat{\bm{G}} = \Reop(\bm{W}\diag(\bm{b})\bm{C})$ matches code exactly. \\
Error denominator & \checkmark & & $\norm{\Sigmat_r}_\Fro^2 = \norm{\mathcal{D}'}_\Fro^2$ (exact since $r=m{-}1$). \\
\midrule
Eigenvalue notation & $\times$ & & Manuscript uses $\gamma_j$; code/SPDMD literature uses $\mu_j$. Conflicts with SPDMD penalty $\gamma$. \\
\bottomrule
\end{tabular}
\end{center}

\medskip\noindent
\textbf{Bottom line:} The manuscript's theoretical formulation and the code implementation are \textbf{mathematically equivalent}.  The manuscript describes the problem in full space ($n_s$-dimensional), while the code works in reduced space ($r$-dimensional) — a standard and exact simplification when $U_r$ has orthonormal columns.  The only issue is the notation $\gamma_j$ for eigenvalues, which should be changed to $\mu_j$ to avoid confusion with the SPDMD sparsity parameter $\gamma$.

\end{document}
